\documentclass[11pt]{article}

\usepackage[margin=2.5cm]{geometry}
\usepackage{amsmath,amssymb,amsthm}
\usepackage{bm}

\title{Solutions Sheet -- Optimization}
\author{}
\date{}

\begin{document}
\maketitle

\subsection*{Part A -- Fundamental Exercises}

\paragraph{Exercise 1. Convex quadratic and unique minimizer.}

Let
\[
f(\bm{x}) = \tfrac{1}{2}\bm{x}^\top Q \bm{x} + \bm{c}^\top \bm{x} + \alpha,
\quad Q=Q^\top.
\]

\begin{enumerate}
\item A twice differentiable function is convex iff its Hessian is positive semidefinite everywhere.  
Here
\[
\nabla^2 f(\bm{x}) = Q \quad (\text{constant in } \bm{x}),
\]
so
\[
f \text{ convex } \Longleftrightarrow Q \succeq 0 \ (\text{positive semidefinite}).
\]

\item Using matrix calculus,
\[
\nabla f(\bm{x}) = Q\bm{x} + \bm{c},\qquad
\nabla^2 f(\bm{x}) = Q.
\]

\item If $Q\succ 0$ (positive definite), $f$ is strictly convex and has a unique minimizer $\bm{x}^\star$.  
First-order optimality:
\[
\nabla f(\bm{x}^\star) = 0
\quad\Longleftrightarrow\quad
Q\bm{x}^\star + \bm{c} = 0
\quad\Longleftrightarrow\quad
\bm{x}^\star = -Q^{-1}\bm{c}.
\]
Strict convexity guarantees uniqueness.
\end{enumerate}

\bigskip

\paragraph{Exercise 2. Subgradients of simple nonsmooth functions.}

\begin{enumerate}
\item For $f(t)=|t|$, the subgradient $g\in\partial f(t)$ satisfies
\[
|y| \ge |t| + g(y-t) \quad \forall y\in\mathbb{R}.
\]
One checks:
\[
\partial f(t)=
\begin{cases}
\{1\}, & t>0,\\[4pt]
[-1,1], & t=0,\\[4pt]
\{-1\}, & t<0.
\end{cases}
\]

\item For $g(\bm{x})=\|\bm{x}\|_1=\sum_{i=1}^n |x_i|$, we use separability:
\[
\partial g(\bm{x}) = \{\bm{s}\in\mathbb{R}^n : s_i\in\partial |x_i|\ \forall i\},
\]
so
\[
s_i =
\begin{cases}
\operatorname{sign}(x_i), & x_i\neq 0,\\[4pt]
\gamma_i,\ |\gamma_i|\le 1, & x_i = 0.
\end{cases}
\]
\end{enumerate}

\bigskip

\paragraph{Exercise 3. Proximal operators: projection and soft-thresholding.}

\begin{enumerate}
\item $g(x)=\iota_{[a,b]}(x)$. We solve
\[
\operatorname{prox}_{\lambda g}(v)
=
\arg\min_{x\in\mathbb{R}}
\left\{
\iota_{[a,b]}(x) + \tfrac{1}{2\lambda}(x-v)^2
\right\}.
\]
Outside $[a,b]$ the function is $+\infty$, so the minimization is restricted to $x\in[a,b]$:
\[
\operatorname{prox}_{\lambda g}(v)
= \arg\min_{x\in[a,b]} (x-v)^2
= P_{[a,b]}(v)
= \min\{\max\{v,a\},b\}.
\]

\item $h(\bm{x})=\lambda\|\bm{x}\|_1$. The proximal operator is separable:
\[
\operatorname{prox}_{h}(\bm{v})
= \bigl(\operatorname{prox}_{\lambda|\cdot|}(v_i)\bigr)_i,
\]
where each coordinate solves
\[
\min_x \lambda|x| + \tfrac12(x-v)^2.
\]
The solution is the soft-thresholding operator:
\[
\operatorname{prox}_{\lambda|\cdot|}(v)
= \operatorname{sign}(v)\,\max\{|v|-\lambda,0\}.
\]
Hence
\[
(\operatorname{prox}_{h}(\bm{v}))_i
= \operatorname{sign}(v_i)\,\max\{|v_i|-\lambda,0\}.
\]
\end{enumerate}

\bigskip

\paragraph{Exercise 4. One step of gradient descent on a quadratic.}

\[
f(\bm{x})=\tfrac{1}{2}\bm{x}^\top Q \bm{x}-\bm{b}^\top\bm{x},\quad Q\succ0.
\]

\begin{enumerate}
\item Gradient:
\[
\nabla f(\bm{x}) = Q\bm{x}-\bm{b}.
\]
Gradient descent:
\[
\bm{x}_{k+1} = \bm{x}_k - \alpha(Q\bm{x}_k - \bm{b})
= (I-\alpha Q)\bm{x}_k + \alpha\bm{b}.
\]

\item Let $0<\lambda_{\min}\le\lambda_{\max}$ be eigenvalues of $Q$.  
Convergence requires spectral radius $\rho(I-\alpha Q)<1$, i.e.
\[
|1-\alpha\lambda_i|<1 \quad \forall i
\quad\Longleftrightarrow\quad
0<\alpha<\frac{2}{\lambda_{\max}}.
\]
Under this condition, $\bm{x}_k\to Q^{-1}\bm{b}$.
\end{enumerate}

\bigskip

\paragraph{Exercise 5. Composite optimization in 1D.}

\[
F(x)=\tfrac12(x-a)^2+\lambda |x|,\quad \lambda>0.
\]

\begin{enumerate}
\item Both terms are convex; sum is convex.

\item Subgradient:
\[
\partial F(x) = (x-a) + \lambda \partial |x|.
\]
We solve $0\in\partial F(x^\star)$ by cases.

\emph{Case 1:} $x>0$, $\partial|x|=\{1\}$:
\[
0 = x-a+\lambda \Rightarrow x=a-\lambda,
\]
valid if $a-\lambda>0$.

\emph{Case 2:} $x<0$, $\partial|x|=\{-1\}$:
\[
0 = x-a-\lambda \Rightarrow x=a+\lambda,
\]
valid if $a+\lambda<0$.

\emph{Case 3:} $x=0$, $\partial|x|=[-1,1]$:
\[
0\in -a+\lambda[-1,1]
\Longleftrightarrow |a|\le\lambda.
\]

\item Summarizing,
\[
x^\star =
\begin{cases}
a-\lambda, & a>\lambda,\\[4pt]
0, & |a|\le\lambda,\\[4pt]
a+\lambda, & a<-\lambda.
\end{cases}
\]
Equivalently,
\[
x^\star = \operatorname{sign}(a)\,\max\{|a|-\lambda,0\},
\]
the soft-thresholding of $a$.
\end{enumerate}

\bigskip

\paragraph{Exercise 6. Two-asset minimum-variance portfolio (finance).}

Weights $\bm{w}=(w_1,w_2)$ with $w_1+w_2=1$, covariance matrix
\[
\Sigma=
\begin{pmatrix}
\sigma_1^2 & \rho\sigma_1\sigma_2\\
\rho\sigma_1\sigma_2 & \sigma_2^2
\end{pmatrix}.
\]

\begin{enumerate}
\item Set $w_2=1-w_1$, so
\[
\bm{w}=\begin{pmatrix}w_1\\1-w_1\end{pmatrix},
\quad
V(w_1)=\bm{w}^\top \Sigma \bm{w}.
\]
Compute:
\[
V(w_1)
= \sigma_1^2 w_1^2 + 2\rho\sigma_1\sigma_2 w_1(1-w_1) + \sigma_2^2 (1-w_1)^2,
\]
a quadratic in $w_1$.

\item Differentiate:
\[
V'(w_1) = 2(\sigma_1^2 - 2\rho\sigma_1\sigma_2 + \sigma_2^2)w_1
+ 2(\rho\sigma_1\sigma_2 - \sigma_2^2).
\]
Set $V'(w_1)=0$:
\[
w_1^\star
=
\frac{\sigma_2^2 - \rho\sigma_1\sigma_2}
{\sigma_1^2 + \sigma_2^2 - 2\rho\sigma_1\sigma_2},
\qquad
w_2^\star=1-w_1^\star
=
\frac{\sigma_1^2 - \rho\sigma_1\sigma_2}
{\sigma_1^2 + \sigma_2^2 - 2\rho\sigma_1\sigma_2}.
\]

\item As $\rho\to1$, assets are perfectly positively correlated, so diversification is ineffective; the optimal allocation tends to favor the lower-variance asset.

As $\rho\to-1$, strong negative correlation allows large risk reduction; the optimizer uses both assets more to exploit diversification.
\end{enumerate}

\bigskip

\paragraph{Exercise 7. Projected gradient in one dimension.}

\[
\min_{x\in[0,1]} f(x)=(x-2)^2.
\]

\begin{enumerate}
\item Unconstrained minimizer: $f'(x)=2(x-2)=0\Rightarrow x=2$.

\item Constrained minimizer: project $2$ onto $[0,1]$:
\[
x^\star = P_{[0,1]}(2)=1.
\]

\item $f'(x)=2(x-2)$, so projected gradient:
\[
x_{k+1} = P_{[0,1]}\bigl(x_k - \alpha f'(x_k)\bigr)
= P_{[0,1]}\bigl(x_k - 2\alpha(x_k-2)\bigr).
\]
Geometrically: take an unconstrained gradient step, then “clip’’ the iterate back to $[0,1]$ if it leaves the interval.
\end{enumerate}

\bigskip

\paragraph{Exercise 8. Least squares regression and normal equations.}

We fit $y\approx ax+b$ to points $(x_i,y_i)$.

\begin{enumerate}
\item Define
\[
\bm{\theta}=
\begin{pmatrix}a\\ b\end{pmatrix},\quad
\bm{y}=
\begin{pmatrix}y_1\\ \vdots\\ y_m\end{pmatrix},\quad
X=
\begin{pmatrix}
x_1 & 1\\
\vdots & \vdots\\
x_m & 1
\end{pmatrix}.
\]
Then
\[
X\bm{\theta} =
\begin{pmatrix}
ax_1+b\\ \vdots\\ ax_m+b
\end{pmatrix}
\]
is the vector of predictions, and
\[
J(\bm{\theta})=\frac12\|X\bm{\theta}-\bm{y}\|^2
\]
is the least squares objective.

\item Gradient:
\[
J(\bm{\theta})=\tfrac12(X\bm{\theta}-\bm{y})^\top(X\bm{\theta}-\bm{y}),
\]
so
\[
\nabla J(\bm{\theta}) = X^\top (X\bm{\theta}-\bm{y}).
\]
Hessian:
\[
\nabla^2 J(\bm{\theta}) = X^\top X.
\]
For any $\bm{z}$,
\[
\bm{z}^\top X^\top X \bm{z} = \|X\bm{z}\|^2 \ge 0,
\]
so $X^\top X$ is positive semidefinite and $J$ is convex (strictly convex if $X$ has full column rank).

\item First-order optimality:
\[
\nabla J(\bm{\theta}^\star)=0
\Longleftrightarrow
X^\top(X\bm{\theta}^\star-\bm{y}) = 0
\Longleftrightarrow
X^\top X \bm{\theta}^\star = X^\top \bm{y}.
\]
These are the \emph{normal equations}. If $X^\top X$ is invertible,
\[
\bm{\theta}^\star = (X^\top X)^{-1} X^\top \bm{y}.
\]

\item $\bm{\theta}^\star=(a^\star,b^\star)$ gives the unique line minimizing the sum of squared residuals.  
If all $x_i$ are equal, columns of $X$ are linearly dependent, $X^\top X$ is singular, and there is no unique best line: only a vertical “cloud” of points.
\end{enumerate}

\bigskip
\bigskip

\subsection*{Part B -- Advanced and Finance-Oriented Exercises}

\paragraph{Exercise 9. Subgradient optimality condition.}

Let $f$ be proper, closed, convex.

\begin{enumerate}
\item If $0\in\partial f(\bm{x}^\star)$, then by definition
\[
f(\bm{x}) \ge f(\bm{x}^\star) + 0^\top(\bm{x}-\bm{x}^\star)
= f(\bm{x}^\star) \quad \forall \bm{x},
\]
so $\bm{x}^\star$ is a global minimizer.

\item Conversely, if $\bm{x}^\star$ is a global minimizer, directional derivatives satisfy $f'(\bm{x}^\star;\bm{d})\ge 0$ for all $\bm{d}$.  
The subdifferential can be characterized by
\[
\partial f(\bm{x}^\star)
= \{\bm{g}: \bm{g}^\top\bm{d} \le f'(\bm{x}^\star;\bm{d}) \ \forall \bm{d}\}.
\]
Since $f'(\bm{x}^\star;\bm{d})\ge 0$, $\bm{g}=\bm{0}$ belongs to $\partial f(\bm{x}^\star)$.

\item This generalizes the smooth condition $\nabla f(\bm{x}^\star)=0$.
\end{enumerate}

\bigskip

\paragraph{Exercise 10. KKT for mean--variance portfolio (finance).}

\[
\min_{\bm{w}} \tfrac12 \bm{w}^\top\Sigma\bm{w}
\quad\text{s.t.}\quad
\bm{1}^\top\bm{w}=1,\quad
\bm{\mu}^\top\bm{w}\ge\bar{r},\quad
\bm{w}\ge\bm{0}.
\]

\begin{enumerate}
\item $\Sigma\succ0$ implies a convex quadratic objective. The constraints are linear, so feasible set is convex: the problem is convex.

\item Define
\[
g_1(\bm{w}) = \bar{r}-\bm{\mu}^\top\bm{w}\le 0,\quad
g_{2,i}(\bm{w}) = -w_i\le 0.
\]
Lagrangian with multipliers $\lambda$ (budget), $\nu\ge0$ (return), $\bm{\gamma}\ge\bm{0}$ (non-negativity):
\[
\mathcal{L}(\bm{w},\lambda,\nu,\bm{\gamma})
=
\tfrac12\bm{w}^\top\Sigma\bm{w}
+ \lambda(\bm{1}^\top\bm{w}-1)
+ \nu(\bar{r}-\bm{\mu}^\top\bm{w})
- \bm{\gamma}^\top\bm{w}.
\]

\item KKT conditions:

Primal feasibility:
\[
\bm{1}^\top\bm{w}^\star=1,\quad
\bm{\mu}^\top\bm{w}^\star\ge\bar{r},\quad
\bm{w}^\star\ge\bm{0}.
\]

Dual feasibility:
\[
\nu^\star\ge0,\quad \bm{\gamma}^\star\ge\bm{0}.
\]

Stationarity:
\[
\nabla_{\bm{w}}\mathcal{L}
=
\Sigma\bm{w}^\star + \lambda^\star \bm{1} - \nu^\star\bm{\mu} - \bm{\gamma}^\star
= \bm{0}.
\]

Complementary slackness:
\[
\nu^\star(\bar{r}-\bm{\mu}^\top\bm{w}^\star)=0,\quad
\gamma_i^\star w_i^\star=0\ \forall i.
\]

\item The vector $\Sigma\bm{w}^\star$ (marginal risk) is balanced by $\lambda^\star \bm{1}$ (budget), $\nu^\star\bm{\mu}$ (target return), and $\bm{\gamma}^\star$ (non-negativity).  
Active constraints (e.g., assets with $w_i^\star=0$) have $\gamma_i^\star>0$ and affect the stationarity condition.
\end{enumerate}

\bigskip

\paragraph{Exercise 11. Dual ascent for equality-constrained quadratic.}

\[
\min_{\bm{x}} \tfrac12\bm{x}^\top Q\bm{x}
\quad\text{s.t.}\quad A\bm{x}=\bm{b},\quad Q\succ 0.
\]

\begin{enumerate}
\item Lagrangian:
\[
\mathcal{L}(\bm{x},\bm{\lambda})
=
\tfrac12\bm{x}^\top Q\bm{x} + \bm{\lambda}^\top(A\bm{x}-\bm{b}).
\]

\item For fixed $\bm{\lambda}$, minimize over $\bm{x}$:
\[
\nabla_{\bm{x}}\mathcal{L} = Q\bm{x} + A^\top\bm{\lambda} = 0
\Rightarrow
\bm{x}(\bm{\lambda}) = -Q^{-1}A^\top\bm{\lambda}.
\]
Plug back:
\[
d(\bm{\lambda})
= \mathcal{L}(\bm{x}(\bm{\lambda}),\bm{\lambda})
= -\tfrac12 \bm{\lambda}^\top A Q^{-1}A^\top \bm{\lambda} - \bm{\lambda}^\top\bm{b},
\]
a concave quadratic.

\item Gradient:
\[
\nabla d(\bm{\lambda})
= -A Q^{-1}A^\top\bm{\lambda} - \bm{b}
= A\bm{x}(\bm{\lambda}) - \bm{b}.
\]

\item Dual ascent:
\[
\bm{\lambda}_{k+1}
= \bm{\lambda}_k + \alpha_k \nabla d(\bm{\lambda}_k)
= \bm{\lambda}_k + \alpha_k (A\bm{x}(\bm{\lambda}_k) - \bm{b}).
\]
The update moves in the direction of the constraint residual $A\bm{x}-\bm{b}$, gradually enforcing feasibility.
\end{enumerate}

\bigskip

\paragraph{Exercise 12. Augmented Lagrangian and method of multipliers.}

Augmented Lagrangian:
\[
\mathcal{L}_\rho(\bm{x},\bm{\lambda})
=
\tfrac12\bm{x}^\top Q\bm{x}
+ \bm{\lambda}^\top(A\bm{x}-\bm{b})
+ \tfrac{\rho}{2}\|A\bm{x}-\bm{b}\|^2.
\]

\begin{enumerate}
\item Method of multipliers:
\[
\bm{x}_{k+1} = \arg\min_{\bm{x}} \mathcal{L}_\rho(\bm{x},\bm{\lambda}_k),
\qquad
\bm{\lambda}_{k+1} = \bm{\lambda}_k + \rho(A\bm{x}_{k+1}-\bm{b}).
\]

\item Stationarity for $\bm{x}_{k+1}$:
\[
Q\bm{x}_{k+1} + A^\top\bm{\lambda}_k
+ \rho A^\top(A\bm{x}_{k+1}-\bm{b}) = 0.
\]
Thus
\[
(Q+\rho A^\top A)\bm{x}_{k+1}
= -A^\top\bm{\lambda}_k + \rho A^\top\bm{b},
\]
a linear system with SPD matrix $Q+\rho A^\top A$.

\item The quadratic penalty term improves conditioning and penalizes violations, often leading to more stable and faster convergence than pure dual ascent.
\end{enumerate}

\bigskip

\paragraph{Exercise 13. ADMM for a regularized portfolio problem (finance).}

\[
\min_{\bm{w}}
\left\{
\tfrac12(\bm{w}-\bm{w}^{\mathrm{ref}})^\top\Sigma(\bm{w}-\bm{w}^{\mathrm{ref}})
+ \lambda\|\bm{w}\|_1
\right\}
\quad\text{s.t.}\quad \bm{1}^\top\bm{w}=1.
\]

\begin{enumerate}
\item Define
\[
f(\bm{w})
=
\tfrac12(\bm{w}-\bm{w}^{\mathrm{ref}})^\top\Sigma(\bm{w}-\bm{w}^{\mathrm{ref}})
+ \iota_{\{\bm{1}^\top\bm{w}=1\}}(\bm{w}),
\qquad
g(\bm{z}) = \lambda\|\bm{z}\|_1.
\]
Then
\[
\min_{\bm{w}} f(\bm{w})+\lambda\|\bm{w}\|_1
\ \text{s.t. }\bm{1}^\top\bm{w}=1
\ \Longleftrightarrow\ 
\min_{\bm{w},\bm{z}} f(\bm{w})+g(\bm{z})
\ \text{s.t. }\bm{w}=\bm{z}.
\]
The splitting isolates:
\begin{itemize}
\item a quadratic + linear constraint part,
\item an $\ell_1$ term with simple proximal operator.
\end{itemize}

\item Scaled ADMM iterations (for some $\rho>0$):

\emph{$\bm{w}$-update:}
\[
\bm{w}^{k+1}
=
\arg\min_{\bm{w}}
\left\{
\tfrac12(\bm{w}-\bm{w}^{\mathrm{ref}})^\top\Sigma(\bm{w}-\bm{w}^{\mathrm{ref}})
+ \iota_{\{\bm{1}^\top\bm{w}=1\}}(\bm{w})
+ \tfrac{\rho}{2}\|\bm{w}-\bm{z}^k+\bm{u}^k\|^2
\right\}.
\]

\emph{$\bm{z}$-update:}
\[
\bm{z}^{k+1}
=
\arg\min_{\bm{z}}
\left\{
\lambda\|\bm{z}\|_1
+ \tfrac{\rho}{2}\|\bm{w}^{k+1}-\bm{z}+\bm{u}^k\|^2
\right\},
\]
which is a soft-thresholding of $\bm{w}^{k+1}+\bm{u}^k$.

\emph{Dual (scaled) update:}
\[
\bm{u}^{k+1}
=
\bm{u}^k + \bm{w}^{k+1} - \bm{z}^{k+1}.
\]

\item Advantage: the $\bm{w}$-step solves a structured quadratic program with one linear equality, while the $\bm{z}$-step is cheap and separable (soft-thresholding).  
This decomposition is efficient and scalable for large $n$ and can be parallelized.
\end{enumerate}

\bigskip

\paragraph{Exercise 14. Transaction costs and nonsmooth portfolio optimization (finance).}

\[
\min_{\bm{w}}
\left\{
\tfrac12\bm{w}^\top\Sigma\bm{w}
+ \lambda\|\bm{w}-\bm{w}^{\mathrm{old}}\|_1
\right\}
\quad\text{s.t.}\quad
\bm{1}^\top\bm{w}=1.
\]

\begin{enumerate}
\item The term $|w_i-w_i^{\mathrm{old}}|$ is the absolute traded amount in asset $i$ (buy or sell).  
Summing and multiplying by $\lambda$ gives total linear transaction costs: $\lambda\|\bm{w}-\bm{w}^{\mathrm{old}}\|_1$.

\item The quadratic $\tfrac12\bm{w}^\top\Sigma\bm{w}$ is convex since $\Sigma\succeq 0$.  
The $\ell_1$-norm is convex, and the sum is convex.  
The constraint $\bm{1}^\top\bm{w}=1$ is affine; thus the problem is convex.

\item Lagrangian with multiplier $\eta$:
\[
\mathcal{L}(\bm{w},\eta)
=
\tfrac12\bm{w}^\top\Sigma\bm{w}
+ \lambda\|\bm{w}-\bm{w}^{\mathrm{old}}\|_1
+ \eta(\bm{1}^\top\bm{w}-1).
\]

\item Subgradient optimality:
\[
\bm{0}
\in
\Sigma\bm{w}^\star
+ \lambda\partial\|\bm{w}^\star-\bm{w}^{\mathrm{old}}\|_1
+ \eta^\star\bm{1},
\qquad
\bm{1}^\top\bm{w}^\star=1.
\]
Let $\bm{s}\in\partial\|\bm{w}^\star-\bm{w}^{\mathrm{old}}\|_1$ with
\[
s_i =
\begin{cases}
\operatorname{sign}(w_i^\star-w_i^{\mathrm{old}}), & w_i^\star\neq w_i^{\mathrm{old}},\\
\gamma_i,\ |\gamma_i|\le1, & w_i^\star = w_i^{\mathrm{old}}.
\end{cases}
\]
Then stationarity is
\[
\Sigma\bm{w}^\star + \lambda\bm{s} + \eta^\star\bm{1} = \bm{0}.
\]

\item The $\ell_1$ term encourages many components of $\bm{w}^\star-\bm{w}^{\mathrm{old}}$ to be exactly zero (subgradient interval $[-1,1]$ at zero).  
This leads to \emph{sparse trades}: only a subset of assets is actually traded.  
Classical mean–variance (without transaction costs) would typically adjust many weights, even for small changes in input parameters.
\end{enumerate}

\bigskip

\paragraph{Exercise 15. Comparing gradient descent and subgradient methods.}

\begin{enumerate}
\item For $f(x)=\tfrac12x^2$, $\nabla f(x)=x$, so gradient descent:
\[
x_{k+1}=x_k-\alpha x_k=(1-\alpha)x_k.
\]
Thus
\[
x_k = (1-\alpha)^k x_0.
\]
If $0<\alpha<2$, then $|1-\alpha|<1$ and $x_k\to 0$ geometrically.

\item For $g(x)=|x|$, subgradient $s_k\in\partial g(x_k)$:
\[
\partial g(x) =
\begin{cases}
\{1\}, & x>0,\\
[-1,1], & x=0,\\
\{-1\}, & x<0.
\end{cases}
\]
Subgradient iteration with $\alpha_k = c/\sqrt{k+1}$:
\[
x_{k+1} = x_k - \alpha_k s_k.
\]
There is no simple closed form for $x_k$; iterates may oscillate near $0$.

\item Qualitatively, gradient descent on $f$ gives monotone, fast (linear) convergence to $0$ in function value and in $x_k$, while for $g$ the subgradient method converges more slowly, with possible oscillations.

\item Theory: for smooth strongly convex functions, gradient descent attains a \emph{linear} rate $f(x_k)-f(x^\star)=O(\rho^k)$; for general nonsmooth convex functions, subgradient methods achieve at best \emph{sublinear} rate $O(1/\sqrt{k})$ in function value with suitable stepsizes.
\end{enumerate}

\end{document}